% !TeX root = ../main.tex
% ====================
% VII. KET LUAN
% ====================
\section{Kết luận}

Bài báo đã nghiên cứu hướng cải tiến TransReID bằng cách bổ sung thông tin ở
mức patch token nhưng vẫn giữ global branch và JPM branch làm đường biểu diễn
chính. Hai thành phần được khảo sát là semantic alignment và local reliability.
Semantic alignment sử dụng human parsing mask như tín hiệu giám sát trong huấn
luyện, giúp patch token học cấu trúc part-level mà không cần mask ở pha suy
luận. Local reliability học độ tin cậy của token cục bộ và chỉ hiệu chỉnh đặc
trưng cuối với hệ số nhỏ nhằm tránh làm suy yếu biểu diễn global.

Trong thiết lập so sánh công bằng trên Market-1501, semantic alignment đạt
$89.5\%$ mAP và $95.3\%$ Rank-1, cao hơn baseline $0.3$ điểm mAP với độ trễ
gần như không đổi. Local reliability đạt $89.9\%$ mAP và $95.2\%$ Rank-1, tăng
$0.7$ điểm mAP nhưng làm độ trễ tăng khoảng $3.45\%$. Bản kết hợp semantic
alignment và local reliability đạt kết quả tốt nhất với $90.1\%$ mAP và
$95.4\%$ Rank-1, cao hơn baseline $0.9$ điểm mAP và $0.3$ điểm Rank-1, trong
khi độ trễ vẫn dưới $5$ ms mỗi ảnh trong điều kiện đo hiện tại.

Các kết quả này cho thấy semantic supervision và reliability-guided local
correction là hai tín hiệu bổ sung hữu ích cho TransReID khi được tích hợp theo
hướng bảo thủ. Trong tương lai, cần đánh giá thêm trên các benchmark có che
khuất nặng như Occluded-Duke, thực hiện thí nghiệm đa seed và khảo sát chất
lượng semantic mask để xác định rõ hơn khả năng tổng quát hóa của phương pháp.
